% Generated draft: requires mathematical and source review.
\documentclass[11pt]{article}
\usepackage[margin=25mm]{geometry}
\usepackage{fontspec}
\setmainfont{Latin Modern Roman}
\usepackage{amsmath,amssymb,mathtools,unicode-math}
\setmathfont{Latin Modern Math}
\usepackage{enumitem,needspace,xurl,hyperref}
\setlength{\parindent}{0pt}
\setlength{\parskip}{5pt}
\setlength{\emergencystretch}{3em}
\newcommand{\N}{\mathbb{N}}
\newcommand{\Z}{\mathbb{Z}}
\newcommand{\Q}{\mathbb{Q}}
\newcommand{\R}{\mathbb{R}}
\newcommand{\C}{\mathbb{C}}
\newcommand{\F}{\mathbb{F}}
\newcommand{\A}{\mathbb{A}}
\newcommand{\PP}{\mathbb{P}}
\newcommand{\E}{\mathbb{E}}
\newcommand{\eps}{\varepsilon}
\newcommand{\dd}{\,\mathrm{d}}
\begin{document}
\section*{UnsolvedMath ID: 20000406}
\textbf{Source problem code:} catalogue entry
\subsection{Definitions and mathematical statement}
The source asks what role wild ramification plays in arithmetic-statistical questions of the kind associated with Melanie Matchett Wood, giving hyperelliptic curves in characteristic two as an example. Wild ramification means ramification with index divisible by the residue characteristic in the curve-cover setting. No particular family measure, statistic or asymptotic regime is specified, so the original question is a research direction.

The supplied report treats separable Artin--Schreier double covers \(y^2+y=f(x)\) in characteristic two. In a reduced presentation let the \(r\) geometric poles have odd jumps \(m_i=2e_i-1\), with \(e_i\ge1\). Its reported specialization of the Pries--Zhu component formula is
\[
 g=\sum_{i=1}^r e_i-1,\qquad
 \dim=g+r-2,\qquad
 \operatorname{codim}=g+1-r=\frac12\sum_{i=1}^r(m_i-1),
\]
and the \(2\)-rank is \(r-1\). Here the \(2\)-rank is the integer \(f_2\) such that the geometric two-torsion of the Jacobian has \(2^{f_2}\) points. The dimension and codimension refer to the geometric hyperelliptic moduli-component interpretation in the report, with varying branch positions, rather than to a fixed equation coefficient space. The customary ambient dimension \(2g-1\) concerns the hyperelliptic moduli space for \(g\ge2\); the source does not provide a complete component statement or low-genus conventions, so those dimension assertions require source review.

A separate candidate exact-distribution theorem fixes \(q=2^s\), \(s\ge1\), and \(k\ge1\), and sets
\[
 d=2k(q-1)+1,\qquad
 f(x)=c+\sum_{\substack{1\le j\le d\\j\text{ odd}}}a_jx^j.
\]
Sample \(c\) and the coefficients \(a_j\) uniformly and independently in \(\F_q\), conditional on \(a_d\ne0\). Let \(C_f\) be the smooth projective curve associated to \(y^2+y=f(x)\). The supplied theorem states that it has genus \(k(q-1)\), \(2\)-rank zero, and exactly the distribution
\[
 \#C_f(\F_q)\ \overset{\mathrm{law}}=\ 1+2B,
 \qquad B\sim\operatorname{Binomial}(q,1/2).
\]
The one denotes the point above infinity. This is a finite-genus coefficient-model assertion, not a limit or an unweighted moduli-space distribution. The report states that conditioning on exact degree is unbiased because \(x^d+x\) belongs to the trace-evaluation kernel and its scalar multiples transport leading coefficients. The kernel here refers to the map taking \(f\) to \((\operatorname{Tr}_{\F_q/\F_2}f(x))_{x\in\F_q}\). Full proofs and bibliographic details are not supplied. These reported results illustrate one wild setting without defining or resolving all statistical questions suggested in the original sentence.
\subsection{Short English statement}
The source asks broadly how wild ramification affects arithmetic statistics. The supplied report describes ramification-stratum dimensions and an exact point-count distribution for a particular uniformly sampled characteristic-two equation family. Its coefficient sampling must be distinguished from counting isomorphism classes of curves.
\subsection{Sources}
\begin{itemize}
\item Local UnsolvedMath registry, permanent ID \(20000406\), source code \texttt{catalogue entry}; AIM algebraic number theory corpus, \texttt{aim-algebraic-number-theory-notes.json}, notes[57]. Attribution: Ulam AI / UnsolvedMath Contributors, CC BY 4.0.
\item AIM workshop \emph{Arithmetic statistics over finite fields and function fields}, Question 28: \url{https://aimath.org/pastworkshops/arithstatffieldproblems.pdf}. The literal attribution is ``Jordan Ellenbird''. The supplied verification regards this as probably a typo for Jordan Ellenberg, citing organizer information and other attributions, but gives no separate workshop-page URL. The literal source name is retained rather than silently corrected.
\item The report names a published Pries--Zhu component formula but supplies no title, date or full citation. Original remarks and literature fields are empty.
\item Registry dated August 14, 2026: statement \texttt{reconstructed\_unverified}, status \texttt{partially\_solved}, classification \texttt{PARTIAL-PROGRESS}; attempt 1 labelled partial result, rigor proved, novelty confidence low. The report is a machine-generated research aid, not peer reviewed, requiring independent verification before citation. Its verification text ends before giving its proposed intended-question formulation.
\end{itemize}
\end{document}
